\chapter{Fundamentos da Manipulação de Dados}
\label{cap:fundamentos_manipulacao_dados}

Se a DDL define a estrutura do banco de dados, a DML (\textit{Data Manipulation Language}) é responsável por operar sobre os dados armazenados nessas estruturas. Em termos práticos, é por meio da DML que um sistema passa a registrar ocorrências do mundo real: alunos são cadastrados, notas são lançadas, matrículas são atualizadas e registros indevidos podem ser removidos.

Esta aula tem como objetivo apresentar os fundamentos da manipulação de dados em SQL, concentrando-se nos comandos \texttt{INSERT}, \texttt{SELECT}, \texttt{UPDATE} e \texttt{DELETE}. Além da sintaxe, o foco recai sobre o raciocínio de uso, erros comuns, boas práticas e situações típicas encontradas em ambientes reais.

\section{A DML no uso cotidiano do banco}

Em um sistema de informação, a estrutura do banco é relativamente estável, mas os dados mudam o tempo todo. Isso significa que a manipulação é muito mais frequente do que a definição estrutural.

Algumas ações típicas do cotidiano correspondem diretamente a comandos DML:

\begin{itemize}
    \item cadastrar um novo aluno;
    \item lançar uma nota;
    \item corrigir o nome de um curso;
    \item alterar o status de uma matrícula;
    \item remover um registro inserido por engano;
    \item consultar alunos ou disciplinas para exibição em tela.
\end{itemize}

Por isso, dominar a DML é essencial para qualquer pessoa que vá desenvolver, administrar ou utilizar sistemas baseados em banco de dados relacionais.

\section{Os quatro comandos fundamentais}

Nesta aula, trabalharemos com quatro comandos principais:

\begin{itemize}
    \item \texttt{INSERT} -- inserir registros;
    \item \texttt{SELECT} -- consultar registros;
    \item \texttt{UPDATE} -- atualizar registros existentes;
    \item \texttt{DELETE} -- remover registros.
\end{itemize}

Embora o \texttt{SELECT} seja frequentemente tratado em um grupo específico de consultas, ele aparece aqui por ser parte indispensável da rotina de manipulação e conferência dos dados.

\section{Banco de exemplo}

Para organizar a prática, adotaremos um banco acadêmico simples, com foco nas tabelas \texttt{alunos} e \texttt{notas}. Esse ambiente é suficientemente pequeno para fins didáticos e, ao mesmo tempo, bastante próximo de cenários reais.

\begin{sqlcode}
CREATE TABLE alunos (
    aluno_id SERIAL PRIMARY KEY,
    nome VARCHAR(100) NOT NULL,
    idade INTEGER,
    curso VARCHAR(100) NOT NULL,
    status VARCHAR(20) DEFAULT 'ATIVO'
);
\end{sqlcode}

\begin{sqlcode}
CREATE TABLE notas (
    nota_id SERIAL PRIMARY KEY,
    aluno_id INTEGER NOT NULL,
    disciplina VARCHAR(100) NOT NULL,
    nota NUMERIC(5,2) NOT NULL,
    data_avaliacao DATE DEFAULT CURRENT_DATE,
    FOREIGN KEY (aluno_id) REFERENCES alunos(aluno_id)
);
\end{sqlcode}

\section{INSERT}

O comando \texttt{INSERT} é usado para adicionar novas linhas a uma tabela.

\subsection{Estrutura geral}

\begin{sqlcode}
INSERT INTO nome_tabela (coluna1, coluna2, coluna3)
VALUES (valor1, valor2, valor3);
\end{sqlcode}

\subsection{Inserção simples}

\begin{sqlcode}
INSERT INTO alunos (nome, idade, curso)
VALUES ('Maria Silva', 22, 'Engenharia de Software');
\end{sqlcode}

Nesse comando:
\begin{itemize}
    \item a tabela destino é \texttt{alunos};
    \item as colunas explicitadas são \texttt{nome}, \texttt{idade} e \texttt{curso};
    \item o identificador \texttt{aluno\_id} será gerado automaticamente;
    \item a coluna \texttt{status} assumirá o valor padrão \texttt{ATIVO}.
\end{itemize}

\subsection{Inserção múltipla}

\begin{sqlcode}
INSERT INTO alunos (nome, idade, curso) VALUES
('Ana Souza', 21, 'Engenharia de Software'),
('Bruno Lima', 24, 'Sistemas de Informação'),
('Carla Ramos', 20, 'Engenharia de Software'),
('Daniel Matos', 23, 'Ciência da Computação');
\end{sqlcode}

A inserção múltipla é útil em cadastros em lote, scripts de carga inicial e testes.

\subsection{Inserção em tabela relacionada}

\begin{sqlcode}
INSERT INTO notas (aluno_id, disciplina, nota)
VALUES (1, 'Banco de Dados', 8.5);
\end{sqlcode}

Nesse caso, o valor informado em \texttt{aluno\_id} precisa corresponder a um aluno já existente na tabela \texttt{alunos}. Caso contrário, a restrição de chave estrangeira impedirá a operação.

\begin{observacao}
A ordem de inserção importa quando existem relações entre tabelas. Primeiro são cadastrados os registros da tabela pai; depois, os da tabela filha.
\end{observacao}

\section{SELECT}

O comando \texttt{SELECT} recupera dados do banco. Mesmo antes de explorar consultas mais avançadas, já é possível realizar operações muito úteis.

\subsection{Selecionando colunas específicas}

\begin{sqlcode}
SELECT nome, idade
FROM alunos;
\end{sqlcode}

\subsection{Selecionando todas as colunas}

\begin{sqlcode}
SELECT *
FROM alunos;
\end{sqlcode}

\subsection{Aplicando filtro com WHERE}

\begin{sqlcode}
SELECT nome, curso
FROM alunos
WHERE curso = 'Engenharia de Software';
\end{sqlcode}

\subsection{Filtrando com operadores relacionais}

\begin{sqlcode}
SELECT nome, idade
FROM alunos
WHERE idade >= 22;
\end{sqlcode}

\subsection{Combinando condições}

\begin{sqlcode}
SELECT nome, curso, status
FROM alunos
WHERE idade >= 21 AND status = 'ATIVO';
\end{sqlcode}

\begin{sqlcode}
SELECT nome, curso
FROM alunos
WHERE curso = 'Engenharia de Software'
   OR curso = 'Sistemas de Informação';
\end{sqlcode}

\section{UPDATE}

O comando \texttt{UPDATE} modifica registros já armazenados.

\subsection{Estrutura geral}

\begin{sqlcode}
UPDATE nome_tabela
SET coluna = valor
WHERE condicao;
\end{sqlcode}

\subsection{Atualização simples}

\begin{sqlcode}
UPDATE alunos
SET idade = 23
WHERE nome = 'Maria Silva';
\end{sqlcode}

\subsection{Atualizando múltiplas colunas}

\begin{sqlcode}
UPDATE alunos
SET idade = 24,
    curso = 'Ciência da Computação'
WHERE aluno_id = 2;
\end{sqlcode}

\subsection{Atualizando vários registros}

\begin{sqlcode}
UPDATE alunos
SET status = 'INATIVO'
WHERE curso = 'Curso Extinto';
\end{sqlcode}

Nesse exemplo, várias linhas podem ser alteradas pela mesma condição.

\begin{danger}
Um \texttt{UPDATE} sem cláusula \texttt{WHERE} pode alterar todas as linhas da tabela. Antes de atualizar, é prudente executar um \texttt{SELECT} com a mesma condição para conferir o conjunto de registros afetados.
\end{danger}

\section{DELETE}

O comando \texttt{DELETE} remove linhas de uma tabela.

\subsection{Exclusão simples}

\begin{sqlcode}
DELETE FROM alunos
WHERE nome = 'Maria Silva';
\end{sqlcode}

\subsection{Exclusão por identificador}

\begin{sqlcode}
DELETE FROM notas
WHERE nota_id = 4;
\end{sqlcode}

\subsection{Exclusão em lote}

\begin{sqlcode}
DELETE FROM notas
WHERE disciplina = 'Tópicos Obsoletos';
\end{sqlcode}

\begin{danger}
Assim como no \texttt{UPDATE}, a ausência de \texttt{WHERE} em um \texttt{DELETE} pode remover todos os registros da tabela.
\end{danger}

\section{Exemplo resolvido 1 -- Cadastro inicial de alunos}

Suponha que a secretaria acadêmica precise iniciar o cadastro de uma nova turma. O primeiro passo é inserir os alunos.

\begin{sqlcode}
INSERT INTO alunos (nome, idade, curso) VALUES
('João Pereira', 19, 'Sistemas de Informação'),
('Lívia Matos', 22, 'Engenharia de Software'),
('Pedro Alves', 20, 'Ciência da Computação'),
('Fernanda Rocha', 23, 'Engenharia de Software');
\end{sqlcode}

Após a inserção, é recomendável verificar o resultado.

\begin{sqlcode}
SELECT *
FROM alunos;
\end{sqlcode}

\textbf{Interpretação:}  
O desenvolvedor confirma se os registros foram inseridos corretamente, se os identificadores foram gerados como esperado e se o valor padrão da coluna \texttt{status} foi aplicado.

\section{Exemplo resolvido 2 -- Lançamento de notas}

Com os alunos cadastrados, o próximo passo é registrar suas notas.

\begin{sqlcode}
INSERT INTO notas (aluno_id, disciplina, nota) VALUES
(1, 'Banco de Dados', 8.5),
(2, 'Algoritmos', 7.0),
(3, 'Estruturas de Dados', 9.0),
(4, 'Banco de Dados', 8.0);
\end{sqlcode}

Em seguida, pode-se verificar o conteúdo da tabela:

\begin{sqlcode}
SELECT *
FROM notas;
\end{sqlcode}

\textbf{Interpretação:}  
Essa etapa confirma se cada nota foi associada ao aluno correto. Também ajuda a identificar erros de digitação em disciplina ou valor numérico.

\section{Exemplo resolvido 3 -- Correção cadastral}

Imagine que o aluno de identificador 3 foi cadastrado no curso errado.

Antes de atualizar, faz-se a conferência:

\begin{sqlcode}
SELECT *
FROM alunos
WHERE aluno_id = 3;
\end{sqlcode}

Depois, realiza-se a correção:

\begin{sqlcode}
UPDATE alunos
SET curso = 'Engenharia de Software'
WHERE aluno_id = 3;
\end{sqlcode}

Por fim, valida-se novamente:

\begin{sqlcode}
SELECT *
FROM alunos
WHERE aluno_id = 3;
\end{sqlcode}

\textbf{Interpretação:}  
Esse fluxo mostra uma boa prática importante: identificar o registro, corrigi-lo e confirmar a alteração.

\section{Exemplo resolvido 4 -- Remoção de nota lançada incorretamente}

Suponha que a nota de identificador 2 tenha sido inserida em disciplina errada.

Primeiro, consulte o registro:

\begin{sqlcode}
SELECT *
FROM notas
WHERE nota_id = 2;
\end{sqlcode}

Depois, remova:

\begin{sqlcode}
DELETE FROM notas
WHERE nota_id = 2;
\end{sqlcode}

Por fim, confirme a exclusão:

\begin{sqlcode}
SELECT *
FROM notas
WHERE nota_id = 2;
\end{sqlcode}

\textbf{Interpretação:}  
Se a consulta final não retornar linhas, a remoção foi concluída com sucesso.

\section{Exemplo resolvido 5 -- Atualização em lote com critério}

Suponha que todos os alunos de um curso descontinuado devam ser movidos para outro.

\begin{sqlcode}
SELECT *
FROM alunos
WHERE curso = 'Análise de Sistemas Antigo';
\end{sqlcode}

\begin{sqlcode}
UPDATE alunos
SET curso = 'Análise e Desenvolvimento de Sistemas'
WHERE curso = 'Análise de Sistemas Antigo';
\end{sqlcode}

\begin{sqlcode}
SELECT *
FROM alunos
WHERE curso = 'Análise e Desenvolvimento de Sistemas';
\end{sqlcode}

\textbf{Interpretação:}  
Esse caso mostra que o \texttt{UPDATE} pode alterar vários registros ao mesmo tempo, desde que o critério seja bem definido.

\section{Erros comuns}

\subsection{Erro comum 1 -- Inserir em tabela filha antes da tabela pai}

Exemplo incorreto:

\begin{sqlcode}
INSERT INTO notas (aluno_id, disciplina, nota)
VALUES (999, 'Banco de Dados', 8.0);
\end{sqlcode}

Se não existir aluno com identificador 999, o banco rejeitará a operação.

\subsection{Erro comum 2 -- Atualizar sem WHERE}

Exemplo perigoso:

\begin{sqlcode}
UPDATE alunos
SET status = 'INATIVO';
\end{sqlcode}

Esse comando altera todos os alunos da tabela.

\subsection{Erro comum 3 -- Deletar sem WHERE}

\begin{sqlcode}
DELETE FROM notas;
\end{sqlcode}

Esse comando remove todas as linhas da tabela \texttt{notas}.

\subsection{Erro comum 4 -- Usar SELECT * sem necessidade}

Embora útil em testes, o uso indiscriminado de \texttt{SELECT *} pode:
\begin{itemize}
    \item dificultar a leitura;
    \item trazer colunas desnecessárias;
    \item acoplar relatórios a estruturas que podem mudar.
\end{itemize}

\section{Boas práticas em DML}

\begin{itemize}
    \item usar \texttt{SELECT} para validar condições antes de \texttt{UPDATE} e \texttt{DELETE};
    \item inserir dados respeitando dependências referenciais;
    \item nomear colunas explicitamente em \texttt{INSERT};
    \item evitar comandos destrutivos sem cópia de segurança;
    \item preferir exclusão lógica quando houver necessidade de histórico.
\end{itemize}

\begin{quadro_dica}
Em muitos sistemas, é preferível mudar o status de um registro para \texttt{INATIVO} em vez de apagá-lo fisicamente. Isso ajuda a preservar histórico e reduz riscos de perda indevida de informação.
\end{quadro_dica}

\section{Minicaso integrador -- Secretaria acadêmica}

Considere a rotina de uma secretaria acadêmica no início do semestre. Ela precisa:

\begin{enumerate}
    \item cadastrar quatro novos alunos;
    \item lançar notas iniciais de avaliação diagnóstica;
    \item corrigir o curso de um aluno inserido incorretamente;
    \item remover uma nota lançada em disciplina errada.
\end{enumerate}

Uma sequência possível seria:

\subsection*{Passo 1 -- Inserção de alunos}

\begin{sqlcode}
INSERT INTO alunos (nome, idade, curso) VALUES
('Aline Castro', 20, 'Engenharia de Software'),
('Rafael Gomes', 21, 'Sistemas de Informação'),
('Bianca Luz', 22, 'Ciência da Computação'),
('Marcelo Reis', 19, 'Engenharia de Software');
\end{sqlcode}

\subsection*{Passo 2 -- Conferência}

\begin{sqlcode}
SELECT *
FROM alunos;
\end{sqlcode}

\subsection*{Passo 3 -- Lançamento de notas}

\begin{sqlcode}
INSERT INTO notas (aluno_id, disciplina, nota) VALUES
(1, 'Banco de Dados', 7.5),
(2, 'Algoritmos', 8.0),
(3, 'Banco de Dados', 6.5),
(4, 'Estruturas de Dados', 9.0);
\end{sqlcode}

\subsection*{Passo 4 -- Correção cadastral}

\begin{sqlcode}
UPDATE alunos
SET curso = 'Sistemas de Informação'
WHERE nome = 'Marcelo Reis';
\end{sqlcode}

\subsection*{Passo 5 -- Remoção de lançamento indevido}

\begin{sqlcode}
DELETE FROM notas
WHERE aluno_id = 3
  AND disciplina = 'Banco de Dados';
\end{sqlcode}

\subsection*{Passo 6 -- Conferência final}

\begin{sqlcode}
SELECT *
FROM alunos;
\end{sqlcode}

\begin{sqlcode}
SELECT *
FROM notas;
\end{sqlcode}

\textbf{Comentário didático:}  
Esse minicaso mostra como os comandos DML aparecem encadeados em um processo de trabalho real. A lógica não está apenas em saber a sintaxe, mas em entender a sequência correta das operações e a necessidade constante de conferência.

\section{Exercícios básicos}

\begin{enumerate}
    \item Insira cinco alunos na tabela \texttt{alunos}.
    \item Cadastre três notas em disciplinas diferentes.
    \item Liste todos os alunos da tabela.
    \item Liste apenas nome e curso.
    \item Atualize a idade de um aluno específico.
    \item Remova uma nota por identificador.
\end{enumerate}

\section{Exercícios intermediários}

\begin{enumerate}
    \item Insira oito alunos distribuídos em ao menos três cursos.
    \item Insira notas para pelo menos metade desses alunos.
    \item Faça uma consulta que retorne apenas alunos com idade maior ou igual a 21.
    \item Atualize o curso de todos os alunos de um curso extinto.
    \item Remova todas as notas de uma disciplina específica.
    \item Escreva a consulta de conferência que você executaria antes de cada \texttt{UPDATE} ou \texttt{DELETE}.
\end{enumerate}

\section{Exercícios avançados}

\begin{enumerate}
    \item Monte um script completo de carga inicial para a tabela \texttt{alunos}.
    \item Monte um script de correção que inclua ao menos dois \texttt{UPDATE}s e um \texttt{DELETE}.
    \item Descreva um cenário em que a exclusão lógica é mais adequada que a exclusão física.
    \item Explique o que pode acontecer se um \texttt{INSERT} tentar referenciar uma chave estrangeira inexistente.
\end{enumerate}

\section{Exercício integrador}

Considere um sistema de biblioteca com as tabelas \texttt{leitor}, \texttt{livro} e \texttt{emprestimo}. Escreva:

\begin{enumerate}
    \item comandos de inserção de leitores;
    \item comandos de inserção de livros;
    \item um comando de registro de empréstimo;
    \item uma consulta de conferência dos empréstimos;
    \item um comando para corrigir a data prevista de devolução;
    \item um comando para remover um empréstimo lançado por engano.
\end{enumerate}

\section{Síntese da aula}

Nesta aula, a Parte III avançou no uso operacional da SQL por meio dos comandos fundamentais da DML: \texttt{INSERT}, \texttt{SELECT}, \texttt{UPDATE} e \texttt{DELETE}. Além da sintaxe, foram apresentados exemplos resolvidos, erros comuns, boas práticas e um minicaso integrador, reforçando a dimensão prática da manipulação de dados. A próxima aula ampliará o \texttt{SELECT} com agregações, subconsultas e junções, elevando o nível analítico das consultas.